\chapter{Achalasia}

\section*{Definition and Overview}
Achalasia is a rare disorder of the oesophagus in which the lower oesophageal sphincter, the muscular valve where the oesophagus meets the stomach, fails to relax properly during swallowing, while the oesophagus simultaneously loses the coordinated wave-like contractions that propel food towards the stomach. As a result, food and liquid accumulate within the oesophagus instead of passing into the stomach. Achalasia is a lifelong condition that develops gradually, usually over years, but treatment is effective in most cases at improving swallowing and restoring quality of life.

\section*{Epidemiology}
Achalasia is uncommon, affecting roughly one to two people in every 100,000 per year, with a lifetime risk of about 1 in 10,000. It can present at any age but most often begins between 30 and 60 years. Men and women are affected about equally, and most cases occur without a family history. Symptoms may be present for months or even years before diagnosis, partly because the early symptoms are easily mistaken for reflux or heartburn.

\section*{Aetiology and Pathophysiology}
The exact cause of achalasia is not fully understood, and in most people no specific cause is identified. The disorder is characterised by the loss of nerve cells in the wall of the oesophagus, particularly in the myenteric plexus, the nerve network that coordinates swallowing. Without these nerves, the lower oesophageal sphincter cannot relax and the oesophageal muscle cannot contract in orderly waves.

\begin{itemize}
    \item \textbf{Immune or inflammatory process:} the nerve damage is thought to result from inflammation, possibly triggered by a viral infection, in people who are genetically susceptible
    \item \textbf{Secondary achalasia:} rarely, the disorder is caused by another condition, such as cancer of the oesophagogastric junction, or by Chagas disease, a parasitic infection that can closely mimic achalasia
\end{itemize}

As swallowing fails, the oesophagus gradually stretches and widens, holding food, fluid, and saliva for long periods.

\section*{Clinical Features}
The symptoms of achalasia develop gradually and are progressive.

\begin{itemize}
    \item Difficulty swallowing (dysphagia) that typically begins with solid foods and later also affects liquids
    \item A sensation that food is stuck or "catching" in the chest
    \item Regurgitation of undigested food, sometimes hours after a meal, particularly at night or when bending over
    \item Coughing or choking at night, and occasionally aspiration of food into the lungs
    \item Chest pain or a squeezing sensation, which may be unrelated to physical activity
    \item Heartburn, which is confusing because achalasia is a failure of the sphincter to open rather than the usual reflux
    \item Gradual weight loss because eating becomes difficult and uncomfortable
\end{itemize}

\section*{Differential Diagnosis}
Several conditions can mimic achalasia and must be excluded.

\begin{itemize}
    \item Gastro-oesophageal reflux disease with a peptic stricture causing narrowing
    \item Oesophageal cancer, which can closely imitate achalasia and is known as pseudoachalasia
    \item Other structural causes of dysphagia, such as an oesophageal web or ring
    \item Other motility disorders, including diffuse oesophageal spasm and nutcracker oesophagus
    \item Eosinophilic oesophagitis, an allergic inflammation of the oesophagus
    \item Scleroderma and other connective tissue diseases affecting the oesophagus
    \item Medication-related oesophageal injury, or swallowing problems from neurological disease
\end{itemize}

\section*{When to Seek Medical Care}
A doctor should be consulted for any swallowing difficulty that is persistent or progressive. It is not normal to need to swallow repeatedly, or to drink large volumes of fluid, to get food down.

\textbf{Call 000 (or local emergency services) for:}
\begin{itemize}
    \item Severe chest pain that could be cardiac in origin
    \item Complete inability to swallow saliva or fluids
    \item Choking or coughing with difficulty breathing
    \item Coughing up blood, or vomit containing blood
    \item Signs of aspiration, such as high fever and breathing difficulty
\end{itemize}

\section*{Diagnosis and Investigations}
Achalasia is confirmed by a combination of imaging and pressure studies.

\begin{itemize}
    \item \textbf{Upper endoscopy (gastroscopy):} a camera examination of the oesophagus that excludes cancer, strictures, and other causes; it is always performed to rule out pseudoachalasia
    \item \textbf{Oesophageal manometry:} the key diagnostic test, which measures pressures within the oesophagus and demonstrates that the lower sphincter does not relax and that peristaltic contractions are absent or disordered
    \item \textbf{Barium swallow:} X-ray imaging during a swallow, which typically shows a narrowed "beak" at the lower end of the oesophagus and a widened, food-filled oesophagus above it
    \item \textbf{CT scan:} may be used to seek an underlying cause in older patients or when the presentation is atypical
    \item \textbf{Timed barium study:} measures how effectively the oesophagus empties after treatment
\end{itemize}

\section*{Management}
The aim of treatment is to weaken or open the lower oesophageal sphincter so that food can pass into the stomach. No treatment restores normal peristalsis, so the focus is on relieving the obstruction at the valve.

\begin{itemize}
    \item \textbf{Pneumatic dilatation:} a balloon is passed to the lower oesophageal sphincter and inflated to stretch it; it is effective but may need to be repeated over the years
    \item \textbf{Laparoscopic Heller myotomy:} keyhole surgery that divides the muscle of the lower sphincter; it is the standard surgical treatment
    \item \textbf{POEM (peroral endoscopic myotomy):} a newer endoscopic technique that divides the sphincter muscle through the mouth without an external incision
    \item \textbf{Botulinum toxin injection:} temporarily relaxes the sphincter and is used mainly in people unable to undergo dilatation or surgery, since the effect lasts only a few months
    \item \textbf{Medicines:} calcium channel blockers or nitrates may provide modest, short-term relief but are rarely used long term
    \item \textbf{After treatment:} heartburn often develops because the sphincter no longer blocks reflux, so acid-suppressing medication is commonly required
    \item \textbf{Follow-up:} regular review and repeat investigations to check how well the oesophagus is emptying
\end{itemize}

\section*{Complications}
\begin{itemize}
    \item Aspiration of food or fluid into the lungs, causing aspiration pneumonia
    \item Weight loss and malnutrition resulting from difficulty eating
    \item Progressive stretching and pouching (dilatation) of the oesophagus over time
    \item Oesophagitis from food and secretions pooling within the oesophagus
    \item Reflux and its complications after treatment, because the sphincter is opened
    \item A small risk of oesophageal cancer after many years of untreated or long-standing disease
    \item Rare complications of treatment, including perforation during dilatation or surgery
\end{itemize}

\section*{Prognosis and Outcomes}
Achalasia is a chronic condition that does not resolve by itself, but treatment is effective for most people: approximately 8--9 out of 10 report significant improvement in swallowing after dilatation or surgery. Symptoms may recur over the years and require repeat treatment, especially after pneumatic dilatation. The risk of serious complications is low when the condition is diagnosed and treated in a specialist centre. People with achalasia need long-term follow-up, including monitoring of the oesophagus, and can generally maintain a good quality of life with treatment.

\section*{Prevention and Self-Care}
Achalasia cannot be prevented, but the following measures help manage symptoms and maintain nutrition.

\begin{itemize}
    \item Eat slowly and chew food thoroughly
    \item Take small meals and drink fluids with meals to help food pass
    \item Sit upright while eating and remain upright for a while afterwards
    \item Avoid very cold or very hot drinks and foods that trigger chest pain or spasm
    \item Avoid eating shortly before bedtime to reduce night-time regurgitation
    \item Raise the head of the bed if night-time cough or regurgitation is troublesome
    \item Maintain nutrition and monitor weight, seeking dietetic help if eating becomes difficult
    \item Attend follow-up appointments and report returning symptoms early, as earlier repeat treatment is easier
\end{itemize}

\section*{Sources and Further Reading}
The following organisations provide reliable information on achalasia for patients, students, and clinicians.

\begin{itemize}
    \item NHS. \textit{Achalasia: symptoms, diagnosis, and treatment.} https://www.nhs.uk/conditions/achalasia/
    \item Mayo Clinic. \textit{Achalasia: overview and treatment options.} https://www.mayoclinic.org/diseases-conditions/achalasia/
    \item MSD Manual. \textit{Achalasia: professional overview.} https://www.msdmanuals.com/
    \item MedlinePlus. \textit{Achalasia health information.} https://medlineplus.gov/ency/article/000267.htm
    \item NICE. \textit{Oesophageal conditions and interventional procedures guidance.} https://www.nice.org.uk/guidance/conditions-and-diseases/
    \item International Foundation for Gastrointestinal Disorders. \textit{Achalasia patient information.} https://aboutgimotility.org/
\end{itemize}
